% Supplementary material for "Hierarchical Transformer Preconditioning for Interactive Physics Simulation"
% Per SIGGRAPH Asia 2026 rules, appendices must be standalone supplementary documents,
% not part of the main paper.tex.
\documentclass[acmtog,anonymous,review]{acmart}

\citestyle{acmauthoryear}
\acmSubmissionID{0000}
\setcopyright{none}
\copyrightyear{2026}
\acmYear{2026}
\acmDOI{10.1145/nnnnnnn.nnnnnnn}
\acmJournal{TOG}
\acmVolume{45}
\acmNumber{6}
\acmArticle{1}
\acmMonth{12}

\usepackage{amsmath}
\usepackage{graphicx}
% Same figure search paths as paper.tex (latexmk from Paper/ or repo root).
\graphicspath{{figures/}{Paper/figures/}}
\usepackage{booktabs}

\begin{document}

\title{Supplementary Material: Hierarchical Transformer Preconditioning for Interactive Physics Simulation}

\author{Anonymous Author(s)}
\affiliation{%
  \institution{Paper under double-blind review}
  \country{}}
\renewcommand{\shortauthors}{Anonymous Author(s)}

\maketitle

This document collects implementation-level details, dataset construction, and additional
tables that support the main paper but cannot fit within the SIGGRAPH Asia 2026 page budget.
Section numbering here is independent of the main paper.
Cross-references from the main PDF use the \texttt{xr-hyper} package with prefix \texttt{SX-};
build this file before \texttt{paper.tex} so \texttt{supplementary.aux} is current.

% ===========================================================================
\section{Full Performance Table with Iteration Counts}
\label{supp:full_perf_table}
% ===========================================================================

Table~2 in the main paper omits PCG iteration counts to fit a single column.
Table~\ref{tab:full_perf_table} reports them here for completeness.

\begin{table}[h]
\centering
\caption{Per-frame mean$\pm$std PCG iterations on \textsf{multiphase\_v2}
($\mathrm{rtol}\!=\!10^{-8}$, $20$ frames per scale).}
\label{tab:full_perf_table}
\footnotesize\setlength{\tabcolsep}{4pt}
\begin{tabular}{@{}l ccccc@{}}
\toprule
Method & $N{=}1024$ & $N{=}2048$ & $N{=}4096$ & $N{=}8192$ & $N{=}16384$ \\
\midrule
Unprec.~CG (GPU)             & $497\!\pm\!144$ & $650\!\pm\!240$ & $1153\!\pm\!434$ & $1765\!\pm\!1744$ & $2002\!\pm\!921$ \\
Jacobi (GPU)                 & $325\!\pm\!83$  & $429\!\pm\!112$ & $839\!\pm\!476$  & $968\!\pm\!261$   & $1480\!\pm\!913$ \\
IC (CPU)                     & $59\!\pm\!18$   & $96\!\pm\!23$   & $113\!\pm\!31$   & $170\!\pm\!47$    & $254\!\pm\!61$ \\
IC / DILU (GPU)              & $11\!\pm\!1$    & $15\!\pm\!1$    & $22\!\pm\!4$     & $30\!\pm\!5$      & $40\!\pm\!3$ \\
AMG (CPU)                    & $5\!\pm\!6$     & $9\!\pm\!6$     & $5\!\pm\!7$      & $7\!\pm\!8$       & $10\!\pm\!8$ \\
AMG / AMGX (GPU)             & $1\!\pm\!0$     & $1\!\pm\!0$     & $1\!\pm\!1$      & $2\!\pm\!0$       & $3\!\pm\!0$ \\
\textbf{Ours (GPU)}          & $\mathbf{50\!\pm\!11}$ & $\mathbf{66\!\pm\!16}$ & $\mathbf{83\!\pm\!24}$ & $\mathbf{174\!\pm\!27}$ & $\mathbf{386\!\pm\!69}$ \\
\bottomrule
\end{tabular}
\end{table}

% ===========================================================================
\section{Multiphase\_v2 Benchmark Generation}
\label{supp:dataset}
% ===========================================================================

This appendix gives the precise procedure for generating each frame of
\textsf{multiphase\_v2}, the benchmark used throughout the main paper. The full generator
is \verb|generate_dataset.py| in the released code (configurable via
\verb|--dataset-type multiphase-v2|).

\paragraph{Grid and ordering.}
Each frame targets $N$ degrees of freedom on a 2D structured grid of dimensions
$W\!\times\!H$ with $W\!\cdot\!H\!\ge\!N$ and $W$, $H$ chosen as close to square as
possible. After construction, all $W\!\cdot\!H$ cells are sorted by Morton (Z-order) code
and truncated to the first $N$ cells, giving a contiguous Morton-ordered index set. Edges
and density values follow the same ordering. We use a structured grid (rather than a
particle cloud) so that the assembled $A$ has the canonical 5-point Laplacian sparsity
pattern, isolating the effect of the heterogeneous coefficient field on conditioning.

\paragraph{Density field.}
Per-frame the heavy density $\rho_{\text{heavy}}$ is drawn log-uniformly from $[5, 100]$
with $\rho_{\text{light}}\!=\!1$. We then sample $n_b\!\sim\!\mathrm{Uniform}\{1,2,3\}$
\emph{barriers}; each barrier is independently
\begin{enumerate}
\item assigned an orientation (vertical or horizontal, equiprobable);
\item given a center coordinate drawn uniformly from $[0.2, 0.8]$ along the cross axis,
in normalized grid coordinates;
\item given a thickness drawn uniformly from $[0.05, 0.20]$;
\item given a gap topology drawn uniformly from $\{\text{top}, \text{bottom},
\text{middle\_hole}, \text{closed}\}$.
\end{enumerate}
A cell is heavy iff it lies in the heavy region of \emph{any} barrier, so multiple
barriers can overlap to form cross- or L-shaped inclusions. Each cell's density is
finally multiplied by $1\!+\!\mathcal{N}(0, 0.05^2)$ for fine-grain symmetry breaking.

\paragraph{Operator assembly.}
$A$ is the standard 5-point pressure-Poisson Laplacian with harmonic-mean face
conductance $w_{ij}=2\rho_i\rho_j/(\rho_i+\rho_j)$, evaluated on cell-cell faces.
Boundary cells use zero-flux (Neumann) boundary conditions, as is standard for
incompressible pressure-projection in graphics-grade fluid simulators. We do \emph{not}
apply a zero-mean constraint to the resulting null-space mode; CG handles it implicitly
when the right-hand side is consistent.

\paragraph{Right-hand side.}
For each frame we draw $\mathbf{b}\!\sim\!\mathcal{N}(0, I)$ projected onto the orthogonal
complement of the constant vector ($\mathbf{1}^\top\mathbf{b}\!=\!0$ to ensure
consistency with the Neumann null-space). This is a worst-case-equivalent right-hand-side
distribution: it weights all spatial frequencies equally, including the low ones that are
hardest for local preconditioners.

\paragraph{Per-scale dataset sizes.}
For each of the five scales $N\in\{1024, 2048, 4096, 8192, 16384\}$, we generate $100$
training frames and $20$ test frames with disjoint seeds. Training and test split the same
generator distribution.

\paragraph{Why this benchmark stresses the architecture.}
Three axes of randomization---contrast, topology, and orientation---combine to make
\textsf{multiphase\_v2} a stress test for any preconditioner with a fixed structural prior.
Closed barriers in particular produce sub-domains coupled \emph{only} through near-zero
conductance, which is the precise regime in which an $\mathcal{H}$-matrix's far-block
representation should out-perform a sparsity-restricted approximate inverse: information
about the pressure jump must travel from one chamber to the other in a single step, a
requirement that local message passing on $\text{nnz}(A)$ cannot satisfy with the depths
typical of practical neural preconditioners.

% ===========================================================================
\section{Out-of-Distribution Generalization}
\label{supp:generalization}
% ===========================================================================

% TODO: Fill in once OOD experiments are complete.
This section will be expanded once the OOD evaluation runs complete; placeholders for
expected sub-sections:
\begin{itemize}
\item OOD scale: train at $N\!=\!8192$, evaluate at $N\!\in\!\{2048, 4096, 16384, 32768\}$
without retraining (currently requires retraining since the partition is keyed to $N$;
this section will explore a fine-tune-only protocol).
\item OOD contrast: train at $\rho_{\text{heavy}}\!\in\![5,100]$, evaluate at
$\rho_{\text{heavy}}\!\in\![100, 1000]$ and at $\rho_{\text{heavy}}\!\in\![1, 5]$.
\item OOD topology: train with $1\!\le\!n_b\!\le\!3$ barriers, evaluate with $n_b\!\ge\!4$
and on the alternative ``fractured-media'' problem class generated by the same
generator script.
\end{itemize}

% ===========================================================================
\section{Precision and Hardware Efficiency}
\label{supp:precision}
% ===========================================================================

% TODO: Fill in once measurements are taken.
This section will report:
\begin{itemize}
\item Sensitivity to apply-path precision (float32, mixed-precision, float64), with PCG
scalar accumulators held at float64 throughout. Expected: float32 stable up to
$N\!\approx\!16384$ at $\mathrm{rtol}\!=\!10^{-8}$; mixed precision needs an outer
refinement loop at the largest scales.
\item Achieved Tensor Core utilization and DRAM bandwidth (via Nsight Compute) for the
inner PCG loop, contrasted with the bandwidth-bound Jacobi and SpMV-bound IC kernels
on the same hardware.
\item Memory footprint of the packed factor tensor as a function of $N$, $L$, and $L_s$.
\end{itemize}

\end{document}
